\chapter{Patterns}
\label{ch:patterns}
\index{patterns}

Patterns describe the shape of a value. They appear in \fc|case| arms, in
\fc|let| definitions, in function parameters and in comprehension
generators. Matching a pattern against a value either fails or succeeds and
binds the pattern's variables.

\begin{center}
\small
\begin{tabularx}{\textwidth}{@{}l X@{}}
\toprule
\textbf{Pattern} & \textbf{Matches} \\
\midrule
\tc{x} & anything, and binds it to \tc{x} \\
\tc{_} & anything \\
\tc{42}, \tc{-1}, \tc{'c'}, \tc{"text"}, \tc{true}, \tc{()} & exactly that literal \\
\tc{[]} & the empty list \\
\tc{p1 :: p2} & a non-empty list whose head matches \tc{p1} and tail \tc{p2} \\
\tc{(p1, p2, p3)} & a tuple, component by component \\
\tc{{ f: p, g }} & a record whose field \tc{f} matches \tc{p}; \tc{g} alone binds field \tc{g} to a variable \tc{g}; other fields are ignored \\
\tc{C} & the constant constructor \tc{C} \\
\tc{C p} & constructor \tc{C} whose payload matches \tc{p} (\tc{Node (l, v, r)}) \\
\tc{(p)} & the same as \tc{p} \\
\bottomrule
\end{tabularx}
\end{center}

Patterns nest freely: \fc|Some (x :: _)| matches an option holding a
non-empty list. A variable may appear only once in a pattern. List
patterns are written with \fc|::| and \fc|[]| (there is no
\fc|[x, y]| pattern: write \fc|x :: y :: []|).

\section{Refutable patterns}
A pattern that can fail --- a literal, a constructor of a type with several
constructors, \fc|[]| or \fc|::| --- is \emph{refutable}. In a \fc|case|,
the next arm is tried. In a \fc|let| or a function parameter there is no
next arm: if the value does not match, evaluation stops with the run-time
error \textit{no pattern matched}. In a comprehension generator, the element
is skipped.

\section{Exhaustiveness and redundancy}
\label{sec:exhaustiveness}
\index{exhaustiveness}\index{warnings}

Both back ends check every \fc|case|, using the type of the value matched:
\begin{itemize}
  \item If some value matches no arm, the \fc|case| is \emph{not
    exhaustive}; the warning shows one such value.
  \item If an arm can never be reached because the arms before it cover
    everything it matches, it is \emph{redundant}.
\end{itemize}
Warnings do not stop the program. Integers, strings, characters and floats
have too many values to enumerate, so a \fc|case| on them needs a variable
or \fc|_| arm to be exhaustive.

\forgeexample{patterns-warnings}

When the missing case does happen, the program stops:

\forgeexample{patterns-runtime-error}

The checker is an implementation of Maranget's algorithm for pattern
matrices~\cite{maranget}; Section~\ref{sec:impl-exhaustiveness} describes it.
